\documentclass[11pt]{article}
\usepackage[margin=0.9in]{geometry}
\usepackage{amsmath,amssymb,mathtools,bm}
\usepackage{booktabs,longtable,tabularx,array}
\usepackage{xcolor}
\usepackage{hyperref}
\usepackage{enumitem}
\usepackage{fancyhdr}
\usepackage{microtype}
\usepackage{caption}
\usepackage{float}
\usepackage{tcolorbox}
\usepackage{siunitx}
\usepackage{listings}
\usepackage{titlesec}
\usepackage{setspace}
\usepackage{lastpage}
\usepackage{tikz}
\usetikzlibrary{arrows.meta,positioning}

\definecolor{navy}{HTML}{17365D}
\definecolor{blue}{HTML}{2F5597}
\definecolor{lightblue}{HTML}{EAF1F8}
\definecolor{green}{HTML}{2F6B4F}
\definecolor{lightgreen}{HTML}{EAF5EF}
\definecolor{orange}{HTML}{B45F06}
\definecolor{lightorange}{HTML}{FFF2CC}
\definecolor{red}{HTML}{9C2F2F}
\definecolor{lightred}{HTML}{FCE8E6}
\definecolor{graybox}{HTML}{F3F4F6}
\definecolor{darkgray}{HTML}{4B5563}

\hypersetup{colorlinks=true,linkcolor=navy,urlcolor=blue,citecolor=navy}
\pagestyle{fancy}
\fancyhf{}
\lhead{KRTI-S-X CGS specification}
\rhead{Initialization and semi-analytic reference solver}
\cfoot{\thepage\ / \pageref{LastPage}}
\setlength{\headheight}{14pt}
\setlength{\parindent}{0pt}
\setlength{\parskip}{0.5em}
\setstretch{1.05}
\titleformat{\section}{\Large\bfseries\color{navy}}{\thesection}{0.6em}{}
\titleformat{\subsection}{\large\bfseries\color{blue}}{\thesubsection}{0.6em}{}
\titleformat{\subsubsection}{\normalsize\bfseries\color{darkgray}}{\thesubsubsection}{0.6em}{}

\newtcolorbox{keybox}[1][]{colback=lightblue,colframe=blue,boxrule=0.8pt,arc=2mm,left=2mm,right=2mm,top=1.5mm,bottom=1.5mm,#1}
\newtcolorbox{warningbox}[1][]{colback=lightorange,colframe=orange,boxrule=0.8pt,arc=2mm,left=2mm,right=2mm,top=1.5mm,bottom=1.5mm,#1}
\newtcolorbox{riskbox}[1][]{colback=lightred,colframe=red,boxrule=0.8pt,arc=2mm,left=2mm,right=2mm,top=1.5mm,bottom=1.5mm,#1}
\newtcolorbox{successbox}[1][]{colback=lightgreen,colframe=green,boxrule=0.8pt,arc=2mm,left=2mm,right=2mm,top=1.5mm,bottom=1.5mm,#1}

\newcommand{\vect}[1]{\bm{#1}}
\newcommand{\Oh}{\hat{\Omega}}
\newcommand{\dd}{\mathrm{d}}
\newcommand{\Ihat}{\widehat I}
\newcommand{\Jhat}{\widehat J}
\newcommand{\Fhat}{\widehat{\vect F}}
\newcommand{\whatv}{\widehat w}
\newcommand{\uhat}{\widehat u}
\newcommand{\phat}{\widehat p}

\lstset{basicstyle=\ttfamily\small,frame=single,breaklines=true,columns=fullflexible,showstringspaces=false}

\begin{document}
\begin{titlepage}
\thispagestyle{empty}
\vspace*{0.7cm}
{\color{navy}\rule{\textwidth}{1.6pt}}\\[0.6cm]
{\Huge\bfseries\color{navy} KRTI-S-X}\\[0.2cm]
{\LARGE\bfseries Frozen CGS Initial/Boundary Conditions}\\[0.2cm]
{\Large and Semi-Analytic Eigenmode Reference Solver}\\[0.7cm]

\begin{tcolorbox}[colback=graybox,colframe=darkgray,boxrule=0.8pt,arc=2mm]
\textbf{Purpose.} This document turns the earlier KRTI-S concept into one concrete, dimensional test. Every physical input needed by the canonical transition-regime case is stated in CGS units. The companion Python program solves the reduced one-dimensional transport/eigenvalue problem and exports the background field and eigenfunctions needed to initialize the growing mode directly.
\end{tcolorbox}

\vfill
\begin{tabularx}{\textwidth}{@{}lX@{}}
\textbf{Frozen case} & KRTI-S-X (transition regime, $\Theta=1$) \\
\textbf{Hydrodynamic reference model} & Inviscid, incompressible, two-fluid Euler equations \\
\textbf{Radiation} & Gray, coherent, isotropic scattering; $\kappa_a=0$ \\
\textbf{Primary result} & Unstable growth rate $s$ and complete linear eigenfunction \\
\textbf{Reference method} & Fourier/Laplace separation + characteristic transport + 1-D linear solve + scalar root \\
\textbf{Companion code} & \texttt{krti\_s\_reference.py} \\
\textbf{Important status} & Research reference prototype; convergence is demonstrated, but the benchmark has not yet been independently peer-reviewed or cross-validated. \\
\end{tabularx}
\vfill
{\small Prepared 30 August 2026.}
\end{titlepage}

\tableofcontents
\newpage

\section{What is frozen in this document}
The previous KRTI-S note intentionally left the dimensional scale open. That is useful for theory but not sufficient for a coding agent. This document freezes a single canonical case, \textbf{KRTI-S-X}, in CGS units.

\begin{keybox}
\textbf{Canonical interpretation.} If two independent codes use the values and boundary conditions in this document, use the supplied eigenfunction normalization, and solve the same incompressible/static-scatterer equations, they are solving the same mathematical problem. No hidden temperature, EOS, source normalization, or boundary choice is required.
\end{keybox}

The thin (KRTI-S-T) and thick (KRTI-S-D) cases remain parameter variants. The transition case is frozen first because it is the scientifically interesting one and is much cheaper to converge than the very optically thick case.

\section{The entire KRTI-S-X input deck in CGS units}

\subsection{Domain and coordinates}
Use the two-dimensional physical domain
\begin{equation}
0\le x<L_x,\qquad -H\le z\le H,
\end{equation}
with
\begin{align}
H &= 1.0000000000000000\ \mathrm{cm},\\
k &= 4.0000000000000000\ \mathrm{cm^{-1}},\\
L_x &= \frac{2\pi}{k}=1.5707963267948966\ \mathrm{cm}.
\end{align}
The unperturbed material interface is at $z=0$. The lower fluid occupies $z<0$ and the upper fluid occupies $z>0$.

There is no physical $y$ dependence in the canonical case. A 3-D code may use a one-cell or finite $y$ extent with periodic $y$ boundaries, but it must preserve $\partial/\partial y=0$.

\subsection{Material and force parameters}
\begin{table}[H]
\centering
\caption{Frozen physical parameters for KRTI-S-X.}
\begin{tabularx}{\textwidth}{@{}lclX@{}}
\toprule
Quantity & Symbol & CGS value & Meaning \\
\midrule
Lower density & $\rho_-$ & $1.0000000000\ \mathrm{g\,cm^{-3}}$ & Light fluid, $z<0$ \\
Upper density & $\rho_+$ & $3.0000000000\ \mathrm{g\,cm^{-3}}$ & Heavy fluid, $z>0$ \\
Atwood number & $A$ & $0.5$ & $(\rho_+-\rho_-)/(\rho_++\rho_-)$ \\
Gravity magnitude & $g$ & $1.0000000000\times10^8\ \mathrm{cm\,s^{-2}}$ & Direction $-\hat z$ \\
Speed of light & $c$ & $2.9979245800\times10^{10}\ \mathrm{cm\,s^{-1}}$ & Physical $c$, no reduced-$c$ in the frozen case \\
Scattering opacity & $\kappa_s$ & $2.0000000000\ \mathrm{cm^2\,g^{-1}}$ & Same specific opacity in both fluids \\
Absorption opacity & $\kappa_a$ & $0$ & No thermal coupling \\
Lower extinction & $\chi_-$ & $2.0000000000\ \mathrm{cm^{-1}}$ & $\kappa_s\rho_-$ \\
Upper extinction & $\chi_+$ & $6.0000000000\ \mathrm{cm^{-1}}$ & $\kappa_s\rho_+$ \\
Lower mean free path & $\lambda_-$ & $0.5\ \mathrm{cm}$ & $1/\chi_-$ \\
Upper mean free path & $\lambda_+$ & $1/6\ \mathrm{cm}$ & $1/\chi_+$ \\
Target radiation support & $\alpha$ & $0.5$ & $\kappa_s F_0/(cg)$ \\
Target net upward flux & $F_0$ & $7.4948114500\times10^{17}\ \mathrm{erg\,cm^{-2}\,s^{-1}}$ & Defines $\alpha=0.5$ \\
Effective downward acceleration & $g_{\rm eff}$ & $5.0000000000\times10^7\ \mathrm{cm\,s^{-2}}$ & $g(1-\alpha)$ \\
\bottomrule
\end{tabularx}
\end{table}

The transition parameter is
\begin{equation}
\Theta=\frac{\kappa_s\rho_{\rm ref}}{k}=1,
\qquad \rho_{\rm ref}=2\ \mathrm{g\,cm^{-3}},
\end{equation}
so the two layers have $\Theta_-=0.5$ and $\Theta_+=1.5$.

\subsection{Frozen bottom radiation intensity}
The exact transport boundary condition is a constant incoming half-space intensity at the bottom. We freeze
\begin{equation}
\boxed{I_b=1.6856350\times10^{18}\ \mathrm{erg\,cm^{-2}\,s^{-1}\,sr^{-1}}.}
\end{equation}
This number is the converged calibration that produces the target net flux $F_0$ for the flat two-layer scattering slab. The corresponding \emph{injected} upward energy flux is
\begin{equation}
F_{\rm in}=\pi I_b=5.29557853\times10^{18}\ \mathrm{erg\,cm^{-2}\,s^{-1}}.
\end{equation}
Most of this injected flux is scattered back out of the lower boundary; the net transmitted vertical flux through the stationary slab is the smaller $F_0$ above.

\begin{warningbox}
The boundary condition is $I_b$, not $F_0$ alone. A transport boundary condition requires the incoming angular intensity. $F_0$ is listed because it determines the hydrostatic support and is an essential equilibrium check.
\end{warningbox}

\subsection{Background gas pressure}
The incompressible equations determine pressure only up to an additive constant. Freeze the interface pressure to
\begin{equation}
\boxed{p_*=p_0(0)=1.0000000000\times10^{12}\ \mathrm{dyn\,cm^{-2}}.}
\end{equation}
The background pressure is
\begin{equation}
p_0(z)=
\begin{cases}
p_*-\rho_-g_{\rm eff}z,&z<0,\\
p_*-\rho_+g_{\rm eff}z,&z>0.
\end{cases}
\end{equation}
Therefore
\begin{align}
p_0(-H)&=1.0000500000\times10^{12}\ \mathrm{dyn\,cm^{-2}},\\
p_0(+H)&=9.9985000000\times10^{11}\ \mathrm{dyn\,cm^{-2}}.
\end{align}
The pressure is continuous at $z=0$ and its gradient changes with density.

\subsection{Initial perturbation amplitude}
Use one horizontal Fourier mode:
\begin{equation}
\eta(x,0)=\eta_0\cos(kx),
\end{equation}
with
\begin{equation}
\boxed{k\eta_0=10^{-3},\qquad \eta_0=2.5000000000\times10^{-4}\ \mathrm{cm}.}
\end{equation}
This is small enough for linear theory and large enough to be measurable in a deterministic calculation. A Monte Carlo calculation may need more histories rather than a larger amplitude if the eigenmode signal is noisy.

\section{Governing equations of the frozen benchmark}

\subsection{Hydrodynamics}
In each fluid,
\begin{align}
\nabla\cdot\vect v&=0,\\
\rho_j\left(\frac{\partial\vect v}{\partial t}+\vect v\cdot\nabla\vect v\right)
&=-\nabla p-\rho_jg\hat{\vect z}+\vect f_r,
\end{align}
where $j\in\{-,+\}$. There is no viscosity, no surface tension, no material energy equation, and no heat conduction.

\subsection{Radiation}
For gray coherent isotropic scattering,
\begin{equation}
\frac{1}{c}\frac{\partial I}{\partial t}+\Oh\cdot\nabla I=\chi(J-I),
\qquad
J=\frac{1}{4\pi}\int I\,\dd\Omega,
\end{equation}
with $\chi=\kappa_s\rho$. The radiation flux and force density are
\begin{equation}
\vect F=\int I\Oh\,\dd\Omega,
\qquad
\boxed{\vect f_r=\frac{\chi}{c}\vect F.}
\end{equation}
The model neglects $O(v/c)$ Doppler and aberration terms. This is part of the frozen mathematical problem, not an accidental numerical approximation.

\section{Every boundary and interface condition}

\subsection{Horizontal hydro boundary}
At $x=0$ and $x=L_x$, all material fields are periodic:
\begin{equation}
q(0,z,t)=q(L_x,z,t)
\end{equation}
for density/material identity, velocity, and pressure.

\subsection{Top and bottom hydro boundaries}
At $z=-1\ \mathrm{cm}$ and $z=+1\ \mathrm{cm}$ use impermeable free-slip walls:
\begin{equation}
\boxed{v_z(x,-H,t)=v_z(x,+H,t)=0.}
\end{equation}
For inviscid flow no tangential-velocity Dirichlet condition is imposed. Numerically, a slip-wall ghost cell should copy the tangential velocity and reverse the normal velocity.

Pressure is \textbf{not} held to a fixed Dirichlet value at these walls. The values $p_0(\pm H)$ above describe the stationary background. A well-balanced ghost construction should extend the hydrostatic profile outside the domain while reflecting $v_z$.

For a ghost-cell center at $z_g<-H$ use
\begin{equation}
\rho_g=\rho_-,\quad p_g=p_*-\rho_-g_{\rm eff}z_g,
\end{equation}
and for $z_g>H$ use
\begin{equation}
\rho_g=\rho_+,\quad p_g=p_*-\rho_+g_{\rm eff}z_g.
\end{equation}
The perturbation pressure is extrapolated consistently with the slip wall; do not clamp it to zero.

\subsection{Horizontal radiation boundary}
Radiation is periodic in $x$. A packet leaving at $x=L_x$ re-enters at $x=0$ with unchanged direction, energy, and time, and vice versa.

\subsection{Bottom radiation boundary, $z=-1\ \mathrm{cm}$}
Let $\mu=\Omega_z$. For incoming directions,
\begin{equation}
\boxed{I(x,-H,\Oh,t)=I_b=1.6856350\times10^{18}\ \mathrm{erg\,cm^{-2}\,s^{-1}\,sr^{-1}},\quad \mu>0.}
\end{equation}
The imposed source is independent of $x$ and $t$. Outgoing directions ($\mu<0$) are not constrained: they leave the domain.

For Monte Carlo injection, the angular probability of packets crossing the boundary is flux weighted:
\begin{equation}
p(\mu)=2\mu,\quad 0<\mu<1,\qquad \phi\sim U(0,2\pi).
\end{equation}
Thus $\mu=\sqrt U$. The packet energy created during a timestep over physical source area $A_b$ must sum to
\begin{equation}
E_{\rm inject}=\pi I_b A_b\Delta t.
\end{equation}

\subsection{Top radiation boundary, $z=+1\ \mathrm{cm}$}
The top is vacuum for incoming directions:
\begin{equation}
\boxed{I(x,+H,\Oh,t)=0,\quad \mu<0.}
\end{equation}
Outgoing directions ($\mu>0$) leave freely.

\subsection{Material interface conditions}
The interface is a material contact with no surface tension and no mass transfer. At the physical interface $z=\eta(x,t)$:
\begin{enumerate}[leftmargin=2em]
\item normal material velocity equals interface velocity;
\item gas pressure is continuous;
\item radiation intensity is continuous for every propagation direction;
\item there is no artificial reflection/refraction law at the material interface;
\item only $\rho$ and therefore $\chi=\kappa_s\rho$ jump.
\end{enumerate}
Tangential velocity may differ between the two inviscid fluids.

\section{The exact initial state}

\subsection{Background fields}
The material background is
\begin{equation}
\rho_0(z)=\begin{cases}1,&z<0,\\3,&z>0,\end{cases}\qquad \vect v_0=0,
\end{equation}
with density in $\mathrm{g\,cm^{-3}}$. The pressure is the $p_0(z)$ profile already given.

The radiation background is \textbf{not isotropic and not spatially constant}. It is the unique stationary solution $I_0(z,\mu)$ of
\begin{equation}
\mu\frac{\partial I_0}{\partial z}+\chi I_0=\chi J_0
\end{equation}
with the exact half-range boundary conditions above. The companion solver exports $I_0(z,\mu)$, $J_0(z)$, and $F_{0z}(z)$.

\begin{warningbox}
Do not initialize the radiation field as $I=I_b$ everywhere, and do not initialize it as an isotropic field chosen from $F_0$. Both are wrong. Use the exported stationary $I_0(z,\mu)$.
\end{warningbox}

\subsection{Eigenmode initialization: preferred method}
The reference solver normalizes the interface eigenfunction to $\widehat\eta=1\ \mathrm{cm}$ and returns complex amplitudes $\widehat q(z)$ under the convention
\begin{equation}
q'(x,z,t)=\Re\left[\eta_0\widehat q(z)e^{st+ikx}\right].
\end{equation}
At $t=0$ initialize:
\begin{align}
\eta(x,0)&=\eta_0\cos(kx),\\
\vect v(x,z,0)&=\eta_0\Re\left[(\uhat(z),\whatv(z))e^{ikx}\right],\\
p(x,z,0)&=p_0(z)+\eta_0\Re\left[\phat(z)e^{ikx}\right],\\
I(x,z,\Oh,0)&=I_0(z,\mu)+\eta_0\Re\left[\Ihat(z,\Oh)e^{ikx}\right].
\end{align}
The density/material identity is assigned from the displaced interface:
\begin{equation}
\rho(x,z,0)=\begin{cases}
\rho_-,&z<\eta_0\cos(kx),\\
\rho_+,&z>\eta_0\cos(kx).
\end{cases}
\end{equation}

\begin{riskbox}
\textbf{Do not shift the background profiles and also add the eigenfunction.} Use $p_0(z)$ and $I_0(z,\mu)$ at the fixed base coordinate $z$, plus the exported perturbations. Replacing $p_0(z)$ by $p_0(z-\eta)$ would double count the first-order interface-displacement term already contained in the eigenmode matching conditions.
\end{riskbox}

\subsection{Linear space-time solution}
Once the unstable root $s$ is known, the entire linear solution is
\begin{equation}
q(x,z,t)=q_0(z)+\eta_0e^{st}\Re\left[\widehat q(z)e^{ikx}\right].
\end{equation}
The interface is
\begin{equation}
\boxed{\eta(x,t)=\eta_0e^{st}\cos(kx).}
\end{equation}
This is what the Python solver means by ``solution as a function of space and time.'' The semi-analytic solver computes $s$ and the one-dimensional eigenfunctions; $x$ and $t$ dependence then follows analytically.

\section{How the semi-analytic solution is actually found}

\subsection{Step 1: stationary one-dimensional radiation field}
For the flat interface the problem is independent of $x$ and $t$:
\begin{equation}
\mu I_0'(z,\mu)+\chi_j I_0=\chi_jJ_0,
\qquad
J_0=\frac12\int_{-1}^{1}I_0\,\dd\mu.
\end{equation}
For constant $J_0$ over a small interval, the transport equation is integrated exactly along a characteristic. The code uses cell-averaged step characteristics. Angular averaging gives a linear system equivalent to a discretized Fredholm equation
\begin{equation}
(\mathcal I-\mathcal K_0)J_0=b_0.
\end{equation}
The solution is reconstructed to give the direction-resolved $I_0$ and net flux.

\subsection{Step 2: Fourier/Laplace ansatz}
Assume the perturbation form
\begin{equation}
e^{st+ikx}.
\end{equation}
Write
\begin{align}
I&=I_0+\Ihat(z,\Oh)e^{st+ikx},\\
\vect v&=(\uhat(z),\whatv(z))e^{st+ikx},\\
p&=p_0+\phat(z)e^{st+ikx},\\
\eta&=\widehat\eta e^{st+ikx}.
\end{align}
The reference normalization is $\widehat\eta=1\ \mathrm{cm}$.

\subsection{Step 3: perturbed angular transport}
Inside either homogeneous layer,
\begin{equation}
\boxed{
\mu\frac{\dd\Ihat_j}{\dd z}+
\left(\chi_j+\frac{s}{c}+ik\xi\right)\Ihat_j
=\chi_j\Jhat_j,
}
\end{equation}
where $\xi=\Omega_x$ and
\begin{equation}
\Jhat=\frac{1}{4\pi}\int\Ihat\,\dd\Omega.
\end{equation}
The imposed incoming boundary source is not perturbed:
\begin{align}
\Ihat(-H,\Oh)&=0,&\mu>0,\\
\Ihat(+H,\Oh)&=0,&\mu<0.
\end{align}

\subsection{Step 4: exact linearized interface source}
Intensity is continuous on the \emph{moving} interface. Expanding at $z=0$ gives
\begin{equation}
\boxed{
\Ihat_+(0,\Oh)-\Ihat_-(0,\Oh)
=-\widehat\eta\left[\partial_z I_0(0,\Oh)\right]_-^+.
}
\end{equation}
Using the background transport equation and continuity of $I_0$ and $J_0$ at the interface,
\begin{equation}
\boxed{
\Ihat_+-\Ihat_-=
-\widehat\eta\,\frac{(\chi_+-\chi_-)}{\mu}\left[J_0(0)-I_0(0,\mu)\right].
}
\end{equation}
This direction-dependent jump is the source that couples the interface corrugation to the radiation field.

\subsection{Step 5: characteristic reduction to a one-dimensional linear system}
For a trial complex $s$, define
\begin{equation}
a_j=\chi_j+s/c+ik\xi.
\end{equation}
Along a characteristic,
\begin{equation}
\Ihat(z)=\Ihat(z_a)e^{-a_j(z-z_a)/\mu}
+\int_{z_a}^{z}\frac{\chi_j}{\mu}\Jhat(z')e^{-a_j(z-z')/\mu}\,\dd z'.
\end{equation}
After angular averaging,
\begin{equation}
\boxed{(\mathcal I-\mathcal K_s)\Jhat=b_s\widehat\eta.}
\end{equation}
This is a one-dimensional complex linear system after quadrature; there is no 2-D radiation grid in the reference calculation. Reconstruct
\begin{equation}
\Fhat(z)=\int\Oh\Ihat\,\dd\Omega
\end{equation}
and then
\begin{equation}
\widehat{\vect f}_{r,j}=\frac{\chi_j}{c}\Fhat_j.
\end{equation}

\subsection{Step 6: forced incompressible hydrodynamics}
Linear incompressibility and momentum give
\begin{align}
ik\uhat+\whatv'&=0,\\
\rho_js\uhat&=-ik\phat+\widehat f_x,\\
\rho_js\whatv&=-\phat'+\widehat f_z.
\end{align}
Eliminating $\uhat$ and $\phat$ gives
\begin{equation}
\boxed{
(\dd_z^2-k^2)\whatv_j
=-\frac{k^2\widehat f_{z,j}+ik\,\partial_z\widehat f_{x,j}}{\rho_js}.
}
\end{equation}
For $\widehat\eta=1\ \mathrm{cm}$,
\begin{align}
\whatv_-(0)&=\whatv_+(0)=s,\\
\whatv_-(-H)&=0,\\
\whatv_+(+H)&=0.
\end{align}
Recover pressure from
\begin{equation}
\boxed{\phat=-\frac{\rho_js}{k^2}\whatv'-\frac{i}{k}\widehat f_x.}
\end{equation}

\subsection{Step 7: scalar dispersion residual}
Pressure continuity on the displaced interface gives
\begin{equation}
\phat_+(0)-\phat_-(0)=(\rho_+-\rho_-)g_{\rm eff}\widehat\eta.
\end{equation}
Define
\begin{equation}
\boxed{D(s)=\phat_+(0)-\phat_-(0)-(\rho_+-\rho_-)g_{\rm eff}\widehat\eta.}
\end{equation}
The unstable reference eigenvalue is
\begin{equation}
\boxed{D(s)=0,\qquad \Re(s)>0.}
\end{equation}
For KRTI-S-X the branch found by the current implementation is real and positive.

\section{What the companion Python program does}
The supplied file \texttt{krti\_s\_reference.py} implements exactly the reduction above using:
\begin{itemize}[leftmargin=2em]
\item Gauss--Legendre quadrature in $\mu=\Omega_z$;
\item uniform azimuthal quadrature in $\phi$;
\item exact cell-averaged step-characteristic transport for a piecewise-constant $J$;
\item dense one-dimensional linear solves for the background and perturbed mean intensity;
\item a one-dimensional finite-difference hydrodynamic BVP after the transport reduction;
\item a safeguarded scalar root solve for $D(s)$.
\end{itemize}

Typical usage is
\begin{lstlisting}
python krti_s_reference.py --case X --nz 128 --nmu 24 --nphi 16 \
    --outdir krti_reference
\end{lstlisting}
Here \texttt{--nz} is the number of radiation cells \emph{per layer}. The program creates a subdirectory \texttt{KRTI-S-X/}.

\subsection{Files intended for initialization}
The most important output is
\begin{lstlisting}
KRTI-S-X_eigenfunctions.npz
\end{lstlisting}
which contains:
\begin{itemize}[leftmargin=2em]
\item $s$, $k$, and $\eta_0$;
\item radiation grid $z$;
\item background $J_0$, $F_{0z}$, $I_0(\mu,z)$;
\item $\Jhat$, $\widehat F_x$, $\widehat F_z$;
\item full direction-resolved $\Ihat(z,\mu,\phi)$;
\item hydrodynamic grid and $\uhat$, $\whatv$, $\phat$;
\item angular quadrature directions and weights.
\end{itemize}
The complex eigenfunctions use the physical-field convention
\begin{equation}
q'(x,z,t)=\Re[\eta_0\widehat q(z)e^{st+ikx}].
\end{equation}

CSV versions are also written for inspection. The compressed NPZ should be preferred for direct code initialization because it preserves complex values and direction-resolved intensity compactly.

\section{Current numerical reference values and convergence status}

\subsection{Background source calibration}
High-resolution one-dimensional background solves give
\begin{equation}
I_b\approx1.685635\times10^{18}\ \mathrm{erg\,cm^{-2}\,s^{-1}\,sr^{-1}}.
\end{equation}
The frozen value above was chosen from this convergence sequence. The physical target net flux is exactly the specified $F_0$; a finite reference discretization should approach it as its $z$ and angular resolution are increased.

\subsection{Unstable-root convergence}
For the frozen $I_b$ and the current reference implementation, representative results are:
\begin{table}[H]
\centering
\caption{Prototype KRTI-S-X root convergence. These values demonstrate that the script is solving a stable branch; they are not yet a peer-reviewed final benchmark table.}
\begin{tabular}{rrrr}
\toprule
Cells/layer & $N_\mu$ & $N_\phi$ & $s\ [\mathrm{s^{-1}}]$ \\
\midrule
64 & 16 & 16 & $1.287158\times10^4$ \\
96 & 20 & 16 & $1.288599\times10^4$ \\
128 & 24 & 16 & $1.289326\times10^4$ \\
160 & 28 & 16 & $1.289754\times10^4$ \\
\bottomrule
\end{tabular}
\end{table}
A simple first-order-in-$\Delta z$ extrapolation places the continuum root near $1.291\times10^4\ \mathrm{s^{-1}}$. A publication-quality reference should add higher-order spatial transport or a more systematic Richardson study before freezing more digits.

For comparison, the classical finite-height RTI result using only the residual effective gravity is
\begin{equation}
s_{\rm cl,eff}=\sqrt{A g_{\rm eff} k\tanh(kH)}=9.99664594\times10^3\ \mathrm{s^{-1}}.
\end{equation}
The current full-transport root is therefore substantially different from the naive effective-gravity value, which is precisely the signal KRTI-S is intended to test.

\subsection{Mandatory analytic self-test}
When $\alpha=0$ the background radiation vanishes and the code must recover
\begin{equation}
s^2=A g k\tanh(kH).
\end{equation}
A coarse internal test ($24$ cells/layer, $N_\mu=N_\phi=8$) recovered the analytic value to approximately $8\times10^{-5}$ relative error. This test should remain in the reference solver as a regression check.

\section{How to initialize a Monte Carlo radiation field}
The exact eigenmode is angularly anisotropic. If a Monte Carlo code can initialize packets from an arbitrary intensity, use
\begin{equation}
I_{\rm init}(x,z,\Oh)=I_0(z,\mu)+\eta_0\Re[\Ihat(z,\Oh)e^{ikx}].
\end{equation}
For the chosen $\eta_0$, the computed intensity perturbation is much smaller than the background (well below one percent in the present reference discretization), so positive-weight packet sampling is not problematic.

A practical procedure is:
\begin{enumerate}[leftmargin=2em]
\item construct a sampling table for $I_{\rm init}$ over cell, $\mu$, and $\phi$ from the NPZ file;
\item sample cell/location proportional to radiation energy density;
\item sample direction from the local discrete/continuous angular distribution;
\item assign packet energies so the ensemble reproduces the integrated $I_{\rm init}$;
\item use many independent transport seeds to separate initialization noise from growth-rate error.
\end{enumerate}

If arbitrary anisotropic packet initialization is unavailable, interface-only initialization is still possible, but early-time transient modes must be discarded before fitting the growth rate.

\section{RICH-compatible low-Mach realization: not exactly the same PDE}
RICH evolves compressible hydrodynamics. The mathematical KRTI-S reference is incompressible. Therefore an ideal-gas RICH calculation can only approximate this benchmark unless an incompressible/barotropic mode is added.

A useful low-Mach proxy is to use $\gamma=5/3$ with the same pressure profile. The specific internal energy is
\begin{equation}
e=\frac{p}{(\gamma-1)\rho}.
\end{equation}
At the unperturbed interface this gives
\begin{align}
e_-(0)&=1.5000000\times10^{12}\ \mathrm{erg\,g^{-1}},\\
e_+(0)&=5.0000000\times10^{11}\ \mathrm{erg\,g^{-1}}.
\end{align}
At the physical bottom and top boundaries,
\begin{align}
e_-(-H)&=1.5000750\times10^{12}\ \mathrm{erg\,g^{-1}},\\
e_+(+H)&=4.9992500\times10^{11}\ \mathrm{erg\,g^{-1}}.
\end{align}
The interface sound speeds are approximately $1.29\times10^6$ and $7.45\times10^5\ \mathrm{cm\,s^{-1}}$, much larger than the RTI velocity scale $\sqrt{g/k}=5.0\times10^3\ \mathrm{cm\,s^{-1}}$.

\begin{warningbox}
A RICH result should be called ``KRTI-S incompressible-limit'' unless convergence with increasing sound speed/background pressure has been demonstrated. The reference Python eigenfunction does not include compressible acoustic corrections.
\end{warningbox}

For a compressible code, run at several $p_*$ values (or sound speeds) and extrapolate the measured growth rate toward the incompressible limit. This is preferable to silently calling one high-pressure ideal-gas calculation exact.

\section{Explicit simulation protocol}
\begin{enumerate}[leftmargin=2em]
\item Build the flat two-layer material state with the exact densities above.
\item Load the stationary $I_0(z,\mu)$ from the deterministic reference package.
\item Load the hydrostatic $p_0(z)$ and set zero background velocity.
\item With no perturbation, verify that the interface does not drift and that the mean radiation force plus gas-pressure gradient balances gravity.
\item Load the eigenfunctions and apply the $\eta_0e^{ikx}$ perturbation to interface, velocity, pressure, and radiation intensity.
\item Evolve with the bottom source, top vacuum, horizontal periodicity, gravity, and scattering momentum exchange all enabled.
\item Measure the fundamental interface Fourier amplitude $A_\eta(t)$.
\item In the linear window fit
\begin{equation}
\ln A_\eta(t)=\ln A_0+s_{\rm num}t.
\end{equation}
\item Compare both $s_{\rm num}$ and selected eigenfunction shapes ($F_x,F_z,w,p$) with the reference package.
\end{enumerate}

A reasonable linear-amplitude upper limit remains
\begin{equation}
A_\eta\lesssim\frac{0.05}{k}=1.25\times10^{-2}\ \mathrm{cm}.
\end{equation}
With $s\simeq1.29\times10^4\ \mathrm{s^{-1}}$, one e-fold is about $7.7\times10^{-5}\ \mathrm{s}$.

\section{What the code must not do}
\begin{itemize}[leftmargin=2em]
\item Do not use absorption/emission or a Fleck factor in KRTI-S-X.
\item Do not replace scattering momentum deposition by a diffusion-pressure-gradient force in the transition case.
\item Do not make $z$ periodic.
\item Do not reflect radiation from the top or bottom; only hydro has slip walls there.
\item Do not initialize radiation as isotropic.
\item Do not perturb only density while calling the state an eigenmode initialization.
\item Do not shift $p_0$ or $I_0$ by the interface displacement and then also add the eigenfunctions.
\item Do not compare a compressible ideal-gas run to the incompressible root without a low-Mach convergence study.
\end{itemize}

\section{Parameter variants after KRTI-S-X}
Keeping all dimensional scales fixed, the earlier thin and thick variants change only $\kappa_s$ (and the bottom source must be recalibrated if the same $\alpha=0.5$ is desired):
\begin{table}[H]
\centering
\begin{tabular}{lrrrr}
\toprule
Case & $\Theta$ & $\kappa_s$ [cm$^2$/g] & $\chi_-$ [cm$^{-1}$] & $\chi_+$ [cm$^{-1}$] \\
\midrule
KRTI-S-T & 0.1 & 0.2 & 0.2 & 0.6 \\
KRTI-S-X & 1 & 2 & 2 & 6 \\
KRTI-S-D & 30 & 60 & 60 & 180 \\
\bottomrule
\end{tabular}
\end{table}
The thick case requires much greater $z$ resolution in the deterministic solver because a uniform physical grid otherwise contains optically thick cells. It should not be treated as numerically converged merely because the linear system solves.

\section{Definition of ``solved'' for the semi-analytic reference}
Before publishing KRTI-S-X reference digits, require:
\begin{enumerate}[leftmargin=2em]
\item background $F_{0z}(z)$ converged and constant;
\item $s$ converged in $N_z,N_\mu,N_\phi$;
\item $\alpha=0$ classical RTI recovered to much better accuracy than the target code tolerance;
\item boundary and interface jump residuals verified explicitly;
\item an independent implementation of the deterministic reference or a second transport method reproduces the same root;
\item eigenfunction normalization and phase conventions frozen in machine-readable files.
\end{enumerate}

\begin{successbox}
\textbf{Bottom line.} The CGS numbers in this document make the physical problem unambiguous. The Python program provides a reproducible \emph{prototype} semi-analytic eigenvalue/eigenfunction solution. The remaining scientific task is not to invent missing initial conditions; it is to converge and independently verify the reference root to publication quality.
\end{successbox}

\newpage
\section*{References}
\addcontentsline{toc}{section}{References}
\begin{thebibliography}{9}
\bibitem{JacquetKrumholz2011}
E. Jacquet and M. R. Krumholz,
``Radiative Rayleigh--Taylor Instabilities,''
\emph{The Astrophysical Journal}, 730:116 (2011).
\href{https://doi.org/10.1088/0004-637X/730/2/116}{doi:10.1088/0004-637X/730/2/116}.

\bibitem{Jiang2013}
Y.-F. Jiang, S. W. Davis, and J. M. Stone,
``Non-linear Evolution of Rayleigh--Taylor Instability in a Radiation Supported Atmosphere,''
\emph{The Astrophysical Journal}, 763:102 (2013).
\href{https://doi.org/10.1088/0004-637X/763/2/102}{doi:10.1088/0004-637X/763/2/102}.

\bibitem{Zhang2026}
C. Zhang et al.,
``Linear stability analysis of radiative effects on compressible Rayleigh--Taylor instability,''
\emph{Physical Review E} 113, 045203 (2026).
\end{thebibliography}

\end{document}
